\section{Data Sources}
\label{sec:data_sources}

We use two publicly available datasets, chosen to satisfy
complementary requirements: scale (many experiment windows for
systematic evaluation) and accessibility (no credentials or
proprietary agreements).

\subsection{REES46 eCommerce Behavior Dataset}
\label{subsec:rees46_source}

The primary new dataset is the REES46 eCommerce behavior
dataset~\citep{rees46dataset}, a large-scale event log from
a multi-category online store.
The data is publicly available on Kaggle and hosted at
\url{https://data.rees46.com}.

\paragraph{Scope.}
The dataset covers 7~months of user activity
(October~2019 through April~2020), comprising approximately
380~million events across an estimated 10+~million unique users.
Each event records a timestamp, event type, product identifier,
category, brand, price, user identifier, and session identifier.

\paragraph{Event types.}
Three event types are recorded:
\begin{itemize}[leftmargin=*]
  \item \texttt{view} --- the user viewed a product page;
  \item \texttt{cart} --- the user added a product to their cart;
  \item \texttt{purchase} --- the user completed a purchase.
\end{itemize}
For our analysis, we define a \emph{trigger} as any event
(view, cart, or purchase) by a user on a given day.
This inclusive definition maximizes the signal available for
modeling user engagement dynamics.

\paragraph{Files.}
The data is distributed as one gzipped CSV per month:

\begin{table}[H]
\centering
\caption{REES46 data files. Sizes are approximate (uncompressed).}
\label{tab:rees46_files}
\small
\begin{tabular}{lrr}
\toprule
\textbf{File} & \textbf{Approx.\ rows} & \textbf{Approx.\ size} \\
\midrule
\texttt{2019-Oct.csv} & 42M & 5.3\,GB \\
\texttt{2019-Nov.csv} & 68M & 8.4\,GB \\
\texttt{2019-Dec.csv} & 67M & 8.0\,GB \\
\texttt{2020-Jan.csv} & 58M & 7.2\,GB \\
\texttt{2020-Feb.csv} & 57M & 7.1\,GB \\
\texttt{2020-Mar.csv} & 50M & 6.0\,GB \\
\texttt{2020-Apr.csv} & 40M & 5.0\,GB \\
\midrule
\textbf{Total} & $\sim$380M & $\sim$47\,GB \\
\bottomrule
\end{tabular}
\end{table}

\paragraph{Why this dataset.}
The REES46 data is uniquely suited to our setting for three reasons.
The original AOAS submission relied on two datasets: a proprietary
collection of 1,774 A/B test arms (which cannot be shared or described
in detail) and the ASOS dataset (76 treatment arms, first-trigger
counts only).
Both limitations were flagged by the reviewers.
REES46 addresses them:
\begin{enumerate}[leftmargin=*]
  \item \textbf{Per-user, per-event timestamps.}
        Unlike the ASOS data used in the original AOAS submission
        (which provides only aggregate first-trigger counts per day),
        REES46 records individual events with timestamps.
        This allows us to construct the full trigger-count matrix
        $A \in \mathbb{N}^{D \times N}$ needed to evaluate the
        \nbmodel{} model.
  \item \textbf{Scale.}
        Seven months of data spanning $\sim$210~days yields
        dozens of non-overlapping 28-day experiment windows,
        enabling systematic calibration studies rather than
        one-off comparisons.
  \item \textbf{Public availability.}
        The data requires no proprietary access, ensuring
        full reproducibility of our analyses.
\end{enumerate}

\subsection{REES46 Data Characteristics}
\label{subsec:rees46_characteristics}
\addresses{R2-minor-a (T2.b), R2-minor-d (T2.c), Editor (T1.d)}

The following visualizations characterize the key features of the
REES46 data that motivate our modeling choices.
These figures directly address Reviewer~2's requests to
``visualize data characteristics'' (R2-minor-a) and to provide
``a reference or a visualization'' for the power-law growth claim
(R2-minor-d).

\begin{figure}[H]
  \centering
  \includegraphics[width=0.85\linewidth]{fig_rees46_cumulative_users.pdf}
  \caption{Cumulative unique users over time for three representative
  28-day experiment windows from the REES46 dataset.\reviewtag{A1}
  The vertical dashed line marks the pilot boundary at $D_0 = 7$ days.
  All three curves exhibit sub-linear (concave) growth, consistent
  with the power-law arrival pattern that the stable-beta prior is
  designed to capture.
  The model must extrapolate the remaining 21~days from only the
  first 7~days of observations.}
  \label{fig:rees46_cumulative}
\end{figure}

\begin{figure}[H]
  \centering
  \includegraphics[width=0.85\linewidth]{fig_rees46_powerlaw_decay.pdf}
  \caption{Log-log scatter of day index versus number of new users
  arriving on that day, for three experiment windows.\reviewtag{A1}
  The approximately linear trend on the log-log scale provides
  direct evidence for the power-law decay of new user arrivals
  \citep{clauset2009power}.
  This pattern is the empirical basis for using stable-beta scaled
  process priors, which produce power-law predictive distributions
  by construction.}
  \label{fig:rees46_powerlaw}
\end{figure}

\begin{figure}[H]
  \centering
  \includegraphics[width=0.85\linewidth]{fig_rees46_new_users_per_day.pdf}
  \caption{Daily new user arrivals across the 28-day experiment window.\reviewtag{A1}
  Day~1 sees $\sim$190K new users; by day~28 this has decayed to
  $\sim$60K.
  The non-stationary decay pattern explains why simple Poisson or
  binomial models (which assume constant arrival rates) are
  inadequate for this data.}
  \label{fig:rees46_daily_new}
\end{figure}

\begin{figure}[H]
  \centering
  \includegraphics[width=0.85\linewidth]{fig_rees46_triggers_per_day.pdf}
  \caption{Daily total trigger volume (all event types) across the
  28-day window.\reviewtag{A1}
  Unlike new user arrivals (which decay), total triggers remain
  relatively stable at 1.2--1.6M per day.
  The gap between total triggers and new users represents re-trigger
  activity from returning users --- the signal that the \nbmodel{}
  model uniquely captures, and that the \bemodel{} and \tgmodel{}
  models cannot represent.}
  \label{fig:rees46_triggers}
\end{figure}

\paragraph{Key statistics.}
Across the 7~non-overlapping 28-day experiment windows:

\begin{figure}[H]
  \centering
  \includegraphics[width=0.8\linewidth]{fig_rees46_triggers_per_user.pdf}
  \caption{REES46: per-user trigger count distribution on a
  log-log scale (15.6M users, 7~months).\reviewtag{A1}
  The linear trend on the log-log scale is direct evidence
  for a power-law distribution of user engagement
  \citep{clauset2009power}.
  24\% of users have exactly 1~trigger; the median is 5;
  the maximum is 199K.}
  \label{fig:rees46_triggers_per_user}
\end{figure}

\begin{figure}[H]
  \centering
  \includegraphics[width=0.8\linewidth]{fig_rees46_days_active.pdf}
  \caption{REES46: number of active days per user
  (15.6M users, 211-day span).\reviewtag{A1}
  49\% of users are active on exactly 1~day;
  74\% on 3~days or fewer.
  This extreme sparsity is the empirical pattern that
  motivates the use of BNP priors with power-law tails.}
  \label{fig:rees46_days_active}
\end{figure}

\noindent
Across the 7~non-overlapping 28-day experiment windows:
pilot sizes range from 957K to 1.8M users ($D_0 = 7$);
total users per window range from 2.8M to 4.4M;
and total triggers per window range from 38M to 63M
(13--18 triggers per user on average).
Approximately 66\% of users in each window first appear
\emph{after} the pilot period, confirming that the prediction
task (forecasting unseen users from pilot data) is both
non-trivial and practically relevant.

\subsection{UCI Online Retail Dataset}
\label{subsec:uci_source}

As a secondary benchmark, we use the UCI Online Retail
dataset~\citep{chen2012data}, a transactional record from a
UK-based non-store online retailer covering
December~2010 through December~2011.
The raw data comprises 541,909 transaction line items across
4,339 unique customers over 374~calendar days.

This dataset is smaller in scale but well-established in the
literature and provides a complementary domain (UK gift retail
vs.\ multi-category eCommerce).
Its inclusion demonstrates that our methodology generalizes
across different data scales and retail contexts.

\paragraph{Key statistics.}
After cleaning (dropping null customer IDs, cancellations,
and negative quantities), 397,924 rows remain across 4,339
customers.
Per-customer trigger counts are heavy-tailed: 36\% of customers
are active on exactly 1~day, 67\% on $\leq 3$~days.
The 13~non-overlapping 28-day experiment windows have pilot
sizes ranging from 34 (low-traffic January) to 474 (holiday
season).

\subsection{ASOS Dataset}
\label{subsec:asos_source}
\addresses{R2-j (T2.a)}

The ASOS dataset~\citep{liu2021datasets} provides first-trigger
counts from 72 online fashion retail experiments (144 treatment
arms: control + treatment per experiment).
Unlike REES46 and UCI, the ASOS data records only the number
of \emph{new} users triggering on each day---not per-user
re-trigger counts.
This means only \bemodel{} and \tgmodel{} can be fitted via
maximum marginal likelihood; \nbmodel{} is fitted via
regression on the cumulative first-trigger curve.

\paragraph{Key statistics.}
All experiments use $\Dpilot = 7$~days of pilot data.
Follow-up lengths $\Dfollow$ vary widely, from 7 to 239~days
(median 55), yielding extrapolation ratios
$\Dfollow / \Dpilot$ from 1.0 to 34.1 (median 7.9).
Pilot sizes range from 25K to 10.5M users
(median 1.4M).

\begin{figure}[H]
  \centering
  \includegraphics[width=0.85\linewidth]{fig_asos_cumulative_users.pdf}
  \caption{ASOS: cumulative new users for three representative
  experiments with short ($D_1 = 10$), medium ($D_1 = 55$),
  and long ($D_1 = 150$) follow-up periods.\reviewtag{A1}
  The vertical dashed line marks the pilot boundary at
  $\Dpilot = 7$.
  All curves show sub-linear growth consistent with the
  power-law arrival pattern.}
  \label{fig:asos_cumulative}
\end{figure}

\begin{figure}[H]
  \centering
  \includegraphics[width=0.75\linewidth]{fig_asos_d1_distribution.pdf}
  \caption{ASOS: distribution of follow-up lengths across
  144 treatment arms.\reviewtag{A1}
  The wide range of $\Dfollow$ (7--239 days) makes ASOS
  a challenging testbed: the model must extrapolate from
  7~pilot days to horizons up to 34$\times$ longer.}
  \label{fig:asos_d1_dist}
\end{figure}

\paragraph{Why this dataset is harder.}
The ASOS experiments differ from REES46 and UCI in three ways
that make prediction harder:
(i)~follow-up horizons are much longer (median ratio 7.9 vs.\ 3.0);
(ii)~only first-trigger data is available, limiting the models
that can be fitted;
and (iii)~fashion retail has more volatile engagement patterns
than general eCommerce.
These factors explain the lower absolute accuracy on ASOS
(\Cref{subsec:asos_stratified}) and make it a useful stress
test for the methodology.

\subsection{Comparison of Data Sources}

\begin{table}[H]
\centering
\caption{Comparison of the three datasets.
ASOS was the only public dataset in the original submission;
REES46 and UCI Online Retail are added in this revision.}
\label{tab:data_comparison}
\small
\begin{tabular}{@{}lccc@{}}
\toprule
& \textbf{ASOS} & \textbf{REES46} & \textbf{UCI Online Retail} \\
\midrule
Source & \citet{liu2021datasets} & Kaggle & UCI ML Repository \\
Granularity & Daily first-triggers & Per-event & Per-transaction \\
Trigger counts & First-trigger only & Full counts & Full counts \\
Users & 25K--10.5M per arm & $\sim$15.6M total & $\sim$4.3K \\
Experiments & 72 (144 arms) & 7+ non-overlapping & 13 \\
$\Dpilot$ & 7 & 7 & 7 \\
$\Dfollow / \Dpilot$ & 1--34 (median 7.9) & 1.8--3.0 & 3.0 \\
Supports \nbmodel{} (MLE) & No & Yes & Yes \\
\bottomrule
\end{tabular}
\end{table}
